
\subsection{电流电阻}

\begin{tabularx}{\columnwidth}{XX}
  电流定义式&$I=\dfrac{q}{t}$\\
  电流微观表达式&$I=nqSv$\\
  电阻定义式&$R=\dfrac{U}{I}$\\
  电阻决定式&$R=\rho\dfrac{l}{S}$\\
  串联电阻&$R=R_1+R_2+\cdots$\\
  并联电阻&$\dfrac1R=\dfrac1{R_1}+\dfrac1{R_2}+\cdots$\\
\end{tabularx}

\begin{formulanotebox}
$R=\dfrac{U}{I}$ 是电阻定义式，不是所有元件都满足线性欧姆定律。灯泡、电动机、二极管等元件要先判断是否可视为纯电阻。

导线理想化时电阻为零，同一根理想导线上的各点可视为等势点。
\end{formulanotebox}

\begin{keypointbox}[title={\strut\textcolor{white}{\bfseries 重点知识点：电路基本认识}}]
\textbf{节点与支路}：节点电流满足流入电流之和等于流出电流之和。复杂电路先找等势点，再判断串并联关系。\\

\textbf{理想电表}：理想电流表内阻为零，可视为导线；理想电压表内阻无穷大，可视为断路。\\

\textbf{稳定状态电容器}：直流稳定后电容器支路断路，但电容器两端仍可能有电压。
\end{keypointbox}

\begin{casetypebox}[title={\strut\textcolor{white}{\bfseries 常见题型：串并联化简}}]
\begin{itemize}[leftmargin=1.5em]
  \item 先把理想导线连接的点合并为同一节点。
  \item 串联电路电流相同，电压按电阻分配。
  \item 并联电路电压相同，电流按电阻反比分配。
  \item 任意部分电阻增大通常会使总电阻增大，反之亦然；含桥式电路时需先判断能否化简。
\end{itemize}
\end{casetypebox}
